\chapter*{ВВЕДЕНИЕ}
\addcontentsline{toc}{chapter}{ВВЕДЕНИЕ}
\label{ch:intro}

В прикладных задачах, таких как робототехника, медицина и финансы, математический вывод динамических систем вида $\dot{x}(t)=f(x(t),u(t))$ крайне затруднен из-за высоких издержек ресурсов физического моделирования. В связи с этим возникает необходимость восстановления нелинейной динамики по дискретным зашумленным данным наблюдений $x_{k}^{obs}=x(t_{k})+\epsilon_{k}$. 

Существующие data-driven подходы по-разному аппроксимируют эволюцию системы и зачастую сравниваются бессистемно, без единого протокола и общих бенчмарков. Это создает пробел в понимании границ применимости конкретных алгоритмов в условиях шума и ограниченных выборок.

\textbf{Цель работы:} провести систематическое теоретико-численное сравнение методов DMD/DMDc, SINDy/SINDy-c и EDMD в унифицированных условиях.

\textbf{Задачи работы:}
\begin{enumerate}
    \item подготовить синтетические наборы данных для LTI с управлением, осциллятора Ван-дер-Поля с управлением и аттрактора Лоренца;
    \item разработать программный пайплайн и применить алгоритмы DMD/DMDc, SINDy/SINDy-c и EDMD с согласованными словарями наблюдаемых;
    \item провести серию экспериментов с вариацией уровня измерительного шума и объема обучающих данных, вычислить метрики качества (MSE, MAE);
    \item выполнить статистический анализ результатов и определить эмпирические границы применимости методов.
\end{enumerate}